\documentclass[10pt,a4paper]{amsart}
\usepackage[margin=19mm]{geometry}
\usepackage{amsmath,amssymb,mathtools}
\usepackage[scr=rsfs,cal=euler]{mathalfa}
\usepackage{xfrac}
\usepackage[unicode,hidelinks,bookmarks=false]{hyperref}
\newcommand{\ie}{i.e.\ }
\newcommand{\cf}{cf.\ }
\newcommand{\resp}{respectively\ }
\newcommand{\Subsection}{\subsection}
\providecommand{\nobreakdash}{\nobreak\hbox{-}}
\setlength{\emergencystretch}{2em}
\multlinegap0pt
\usepackage{bm}
\usepackage{bbm}
\usepackage{dsfont}
\usepackage{xspace}
\usepackage{cancel}
\usepackage{mathtools}
\usepackage{amsthm}
\makeatletter
\@ifundefined{theorem}{\newtheorem{theorem}{Theorem}}{}
\makeatother
\newcommand{\Exp}{\operatorname{Exp}}
\newcommand{\RR}{\mathbb{R}}
\newcommand{\SSS}{\mathbb{S}}
\title[Radial Compensation: The Inverse Base-Distribution Problem f]{Radial Compensation: The Inverse Base-Distribution Problem for Chart-Based Generative Models on Riemannian Manifolds}
\author{Marios Papamichalis, Regina Ruane}
\date{Source: arXiv:2511.14056v3 (CC BY 4.0)}
\begin{document}
\maketitle

\noindent\emph{Source and licence.} Radial Compensation: The Inverse Base-Distribution Problem for Chart-Based Generative Models on Riemannian Manifolds. Version arXiv:2511.14056v3 (CC BY 4.0), distributed under the Creative Commons Attribution 4.0 International licence, as recorded in the arXiv licence metadata for this exact version and shown on the abstract page https://arxiv.org/abs/2511.14056v3. Full rights notice: licenses/arxiv-2511.14056-CC-BY-4.0.txt.\par\smallskip

\noindent\textit{Editorial scope.} Counted unit: the Theorem whose source label is \texttt{thm:balanced-polar} in the article (source0.tex lines 1370--1388, source label \texttt{thm:balanced-polar}), delivered together with the article's own complete original proof of it (source0.tex lines 1390--1412, one \texttt{proof} environment of 23 lines). No mathematics is written, repaired or re-derived: statement and proof are the article's own text line for line. Required context retained uncounted: eq:rc-base (lines 541--547); thm:rc-invariances (lines 551--565). Every definition, display and label of those lines is retained unchanged. Omitted from this package and not redistributed: 1--540, 548--550, 566--1369, 1389--1389, 1413--2738. Nothing inside a retained excerpt is omitted or altered: every retained body is delivered exactly as the article's lines are written, so no omission receipt of any kind accompanies this package. Citations: the retained excerpts carry no citation command at all, so no bibliography entry of the article is transcribed and no reference is redistributed. Reference adaptation: none was needed and none was made; every reference inside the delivered unit resolves inside it, so no unresolved reference or citation is produced. Printed slips kept literal: no printed word, symbol, hypothesis or number of the counted statement or of its proof was altered. Layout and class: the article's own class is \texttt{article} with packages including iclr2027\_conference, times, natbib, iclr2026\_conference, neurips\_2026, geometry, caption, algorithm, algorithmic, amsmath, amssymb, amsfonts, amsthm, mathtools; the wrapper is the lane's pinned \texttt{amsart} editorial wrapper at 10pt on A4, which restates the theorem-like environments the retained text needs and adds the article's own macro definitions that the retained text needs (\texttt{\textbackslash Exp}, \texttt{\textbackslash RR}, \texttt{\textbackslash SSS}), with extra wrapper packages (bm, bbm, dsfont, xspace, cancel, mathtools). The delivered numbering is the unit's own. Assets: no retained span contains a figure environment, a graphics-inclusion command, a PDF-page inclusion, a table or a TikZ picture; no image file is copied into this package, the package declares no asset dependency and no image file is redistributed with it. Licence: version arXiv:2511.14056v3 of this article is distributed under the Creative Commons Attribution 4.0 International licence, as recorded in the arXiv licence metadata for this exact version and shown on the abstract page https://arxiv.org/abs/2511.14056v3. A full rights notice is delivered at licenses/arxiv-2511.14056-CC-BY-4.0.txt. Characters: every retained excerpt is pure ASCII, so the delivered engine, XeLaTeX with its default Latin Modern text font, reports no missing character.
\begin{equation}
  f_\theta(x)
  \;:=\;
  \varphi_\theta\!\bigl(R_T(\|x\|)\bigr)\,J_T(\|x\|)\,\mathbf{1}\{x\in U\},
  \qquad x\in\RR^n.
  \label{eq:rc-base}
\end{equation}

\begin{theorem}[Geodesic-radial Fisher and KL invariance under RC]
\label{thm:rc-invariances}
Let $T$ be a scalar-Jacobian azimuthal chart about $p$, let $X\sim f_\theta$
be the RC base defined by~\eqref{eq:rc-base}, and let $Q=T(X)$.
Then for any parameter $\theta$ that enters only through the radius law
$p_{R,\theta}$ (with $\kappa$ fixed):
\begin{align}
\mathcal{I}_M(\theta) &= \mathcal{I}_{R^1}(\theta), \label{eq:fisher-eq}\\
\mathrm{KL}(\rho_\theta\,\|\,\rho_\eta) &= \mathrm{KL}(p_{R,\theta}\,\|\,p_{R,\eta}).
\label{eq:kl-eq}
\end{align}
The pushforward density on $M$ is geodesic-radial,
$\rho_\theta(q)=\varphi_\theta(d(p,q))$,
under standard dominated-score regularity.
\end{theorem}

\begin{theorem}[Balanced polar pushforward beyond constant curvature]
\label{thm:balanced-polar}
Let $(M,g)$ admit geodesic polar coordinates about $p$ on a star-shaped domain $U$ with
$d\mathrm{vol}_M=J_p(R,\omega)\,dR\,d\omega$.
Let $S_M(R):=\int_{\SSS^{n-1}}J_p(R,\omega)\,d\omega$ and let $\varphi_\theta$ be a 1D density
w.r.t.\ $S_M(R)\,dR$ on $[0,R_{\max})$.
Define a ``balanced polar'' density on $[0,R_{\max})\times\SSS^{n-1}$ (w.r.t.\ $dR\,d\omega$) by
\[
f^{\mathrm{bal}}_\theta(R,\omega) := \varphi_\theta(R)\,J_p(R,\omega).
\]
Let $Q=\Exp_p(R\omega)$ with $(R,\omega)\sim f^{\mathrm{bal}}_\theta$. Then the induced density on $U$
(w.r.t.\ $d\mathrm{vol}_M$) is geodesic-radial:
\[
\rho_\theta(q)=\varphi_\theta\!\bigl(d(p,q)\bigr),
\]
and the radius $R=d(p,Q)$ has density $p_{R,\theta}$ w.r.t.\ Lebesgue $dR$, equivalently $\varphi_\theta$ w.r.t.\ $S_M(R)\,dR$.
Moreover, under the usual dominated-score regularity, the Fisher information in $\theta$ equals
the 1D Fisher functional w.r.t.\ $S_M(R)\,dR$.
\end{theorem}

\begin{proof}
Let $F(R,\omega)=\Exp_p(R\omega)$.
For any bounded measurable $g:U\to\RR$,
\[
\int_U g(q)\,\rho_\theta(q)\,d\mathrm{vol}_M(q)
=
\int_0^{R_{\max}}\!\!\int_{\SSS^{n-1}}
g\!\bigl(F(R,\omega)\bigr)\,\rho_\theta\!\bigl(F(R,\omega)\bigr)\,J_p(R,\omega)\,d\omega\,dR.
\]
By definition of the pushforward of $f^{\mathrm{bal}}_\theta$,
\[
\int_U g(q)\,d\nu_\theta(q)
=
\int_0^{R_{\max}}\!\!\int_{\SSS^{n-1}}
g\!\bigl(F(R,\omega)\bigr)\,\varphi_\theta(R)\,J_p(R,\omega)\,d\omega\,dR.
\]
Equality for all $g$ implies $\rho_\theta(F(R,\omega))=\varphi_\theta(R)$ for a.e.\ $(R,\omega)$,
i.e.\ $\rho_\theta(q)=\varphi_\theta(d(p,q))$ on $U$.

Integrating out $\omega$ yields the radius law w.r.t.\ $S_M(R)\,dR$.
The Fisher reduction is the same as in Theorem~\ref{thm:rc-invariances}:
the score depends on $q$ only through $R$, then integrate in polar coordinates.
\end{proof}

\end{document}
